% ============================================================
% UNIDAD 7: TERMINALES, EMPALMES, CRIMPADO Y SOLDADURA
% ============================================================
\clearpage
\section*{UNIDAD 7: TERMINALES Y CRIMPADO}
\subsection*{Pelado, terminales, crimpado, inspección y soldadura}

\begin{recuadro}
\textbf{Objetivo de la unidad:} pelar cables sin dañarlos, seleccionar terminales y empalmes,
crimpar con la herramienta correcta, inspeccionar y ensayar el resultado, y terminar blindajes y
conexiones soldadas según las prácticas aeronáuticas.
\end{recuadro}

% ------------------------------------------------------------
\section*{1. EL PELADO}

\begin{itemize}
  \item Usar \textbf{peladoras de precisión} con la cuchilla correspondiente a la sección y al tipo
        de aislación (o peladoras térmicas para aislaciones delgadas). Nunca pelar con alicate o
        cuchillo.
  \item No cortar ni marcar los hilos: un hilo cortado reduce la sección y es un punto de rotura
        por vibración. La norma admite muy pocos hilos dañados, según la sección.
  \item No retorcer los hilos más de lo que ya vienen trenzados; no estañar un cable que se va a
        crimpar.
  \item La longitud de pelado la indica el fabricante del terminal o del contacto.
\end{itemize}

% ------------------------------------------------------------
\section*{2. TERMINALES Y EMPALMES}

\begin{center}
\small
\begin{tabular}{|p{0.26\linewidth}|p{0.62\linewidth}|}
\hline
\textbf{Tipo} & \textbf{Uso} \\ \hline
Terminal de ojal (anillo) & El preferido en aeronaves: no se sale aunque se afloje el tornillo.
Aislado (con color según la sección) o desnudo. Familias aprobadas como MS25036 o PIDG. \\ \hline
Terminal de horquilla & Solo donde el fabricante lo permite; puede salirse si se afloja. \\ \hline
Terminales para cables gruesos & Para cables de batería y arranque: se crimpan con herramientas
hidráulicas. \\ \hline
Empalmes a tope (\textit{butt splices}) & Unen dos cables en reparaciones. Se evitan en zonas de
movimiento y se escalonan en un mazo. \\ \hline
Contactos de conector & Pines y casquillos que se crimpan e insertan en conectores (Unidad 8). \\
\hline
Manguitos de soldadura & Terminación de blindajes y uniones selladas. \\ \hline
\end{tabular}
\end{center}

\paragraph*{Código de colores de los terminales aislados:} rojo para 22--18 AWG, azul para 16--14
AWG y amarillo para 12--10 AWG.

% ------------------------------------------------------------
\section*{3. EL CRIMPADO}

\subsubsection*{3.1. Qué es un buen crimpado}
El crimpado deforma el barril del terminal sobre los hilos hasta formar una unión \textbf{hermética
a los gases}: metal con metal, sin huecos por donde entre humedad que corroa. Es más confiable que
la soldadura ante la vibración, porque no crea una zona rígida donde los hilos se quiebran.

\subsubsection*{3.2. La herramienta correcta}
\begin{itemize}
  \item \textbf{Pinzas de crimpar con trinquete:} no se abren hasta completar el ciclo, lo que
        asegura una presión constante. Cada matriz corresponde a una familia de terminales y a un
        rango de secciones (color).
  \item \textbf{Crimpadoras para contactos de conectores} con norma militar: la \textbf{M22520/1-01}
        (conocida como AF8) para contactos de 12 a 26 AWG y la \textbf{M22520/2-01} (AFM8) para
        contactos pequeños (de 20 AWG y menores). Se usan con un \textbf{posicionador} o
        \textbf{torreta} específico para cada contacto y un \textbf{selector} de la sección del
        cable.
  \item \textbf{Calibración:} las herramientas se verifican periódicamente con calibres
        pasa/no-pasa (\textit{go/no-go}). Una herramienta fuera de calibración produce crimpados
        flojos que pasan la inspección visual.
  \item \textbf{Nunca} usar alicates comunes ni crimpadoras automotrices sin trinquete.
\end{itemize}

\subsubsection*{3.3. Procedimiento}
\begin{enumerate}
  \item Verificar el número de parte del terminal o contacto y la sección del cable.
  \item Seleccionar la herramienta, la matriz o el posicionador y el selector de sección.
  \item Pelar a la longitud indicada.
  \item Insertar el cable hasta que los hilos se vean en la ventana de inspección (o lleguen al
        tope), sin que la aislación entre en el barril.
  \item Crimpar completando el ciclo del trinquete.
  \item Inspeccionar y, en los ensayos de calidad, realizar el ensayo de tracción.
\end{enumerate}

\subsubsection*{3.4. Inspección}
\begin{center}
\small
\begin{tabular}{|p{0.44\linewidth}|p{0.44\linewidth}|}
\hline
\textbf{Correcto} & \textbf{Defectos} \\ \hline
Hilos visibles en la ventana de inspección o al ras del extremo del barril. & Hilos que no llegan
o hilos sueltos fuera del barril. \\ \hline
Aislación cerca del barril, sin entrar en él (pequeña separación visible). & Aislación atrapada
dentro del crimpado o conductor desnudo expuesto. \\ \hline
Marca de la matriz centrada y simétrica. & Crimpado descentrado, barril fisurado o deformado. \\
\hline
El cable no gira ni se desplaza dentro del terminal. & El cable gira o sale al tirar. \\ \hline
\end{tabular}
\end{center}

\paragraph*{Ensayo de tracción.} Sobre muestras crimpadas con la misma herramienta y el mismo lote,
se aplica una fuerza de tracción especificada para cada sección (valores en la norma del
terminal o contacto o en la AC 43.13-1B). El crimpado debe resistirla sin que el cable se salga. Es
la forma objetiva de validar una herramienta y un operador.

% ------------------------------------------------------------
\section*{4. SOLDADURA Y MANGUITOS}

\begin{itemize}
  \item La soldadura se usa en los contactos de soldar de algunos conectores, en placas y en las
        terminaciones de blindaje con \textbf{manguitos de soldadura}.
  \item Estaño con fundente de colofonia no corrosivo; estación de soldado con temperatura
        controlada; limpiar el fundente después.
  \item La unión debe quedar brillante, con el estaño ``mojando'' el conductor y el contacto, sin
        bolitas ni puntas, y sin que el estaño suba por los hilos más allá del contacto (eso los
        vuelve rígidos y se quiebran).
  \item \textbf{Manguitos de soldadura} (familia M83519): tubo termocontraíble con un anillo de
        estaño y fundente en el centro y anillos selladores en los extremos. Con la pistola de calor
        se contrae el tubo, se funde el estaño que une la malla con el cable de drenaje y se sella
        la unión. Se inspecciona que el anillo se haya fundido por completo.
  \item \textbf{Termocontraíbles} de identificación y protección: con la relación de contracción
        adecuada; calentar parejo, sin quemar.
\end{itemize}

% ------------------------------------------------------------
\section*{5. ACTIVIDADES DE LA UNIDAD}
\begin{enumerate}
  \item Pelar 10 cables de 22 AWG con una peladora de precisión e inspeccionarlos con lupa.
  \item Crimpar terminales de ojal de los tres colores y empalmes a tope; inspeccionar y hacer el
        ensayo de tracción.
  \item Crimpar contactos de conector con la M22520/2-01 y su posicionador; inspeccionar.
  \item Terminar el blindaje de un cable M27500 con manguito de soldadura y cable de drenaje.
  \item Verificar una crimpadora con el calibre pasa/no-pasa y registrar el resultado.
\end{enumerate}
